% !TEX root = ../main.tex
\section{Problem Definition and Representation}

A vacuum-cleaning robot works in a rectangular room divided into cells. Some cells are obstacles,
some are dirty, and one is the charging station. The robot must (1) choose a time slot for every
cleaning task so that all scheduling constraints hold, then (2) navigate from the charger to each
dirty cell in slot order and back to the charger, using graph search.

\textbf{Environment grid.} \code{Grid} is an immutable record with \code{rows}, \code{cols}, the set
of \code{obstacles}, the dirty cells \code{dirty} (label $\to$ cell, in file order, labels
D1\dots Dn) and the \code{charger} cell (\dec{06}). Coordinates are zero-based \code{(row, col)}
with \code{(0, 0)} at the top left. Dirty cells and the charger are passable; only obstacles and
out-of-bounds cells are blocked.

\textbf{Search problem (one navigation leg).}

\begin{center}
\begin{tabularx}{\linewidth}{@{}lX@{}}
\toprule
Element & Representation \\
\midrule
State & the robot's cell \code{(row, col)} \\
Initial state & the leg's start cell (charger or previous dirty cell) \\
Actions & UP $(-1,0)$, DOWN $(+1,0)$, LEFT $(0,-1)$, RIGHT $(0,+1)$, tried in that order (\dec{06}) \\
Transition model & \code{get\_successors(state, grid)}: apply the action; keep the result if it is inside the grid and not an obstacle \\
Goal test & \code{goal\_test(state, goal\_state)}: the robot is on the leg's target cell \\
Step cost & 1 per move; \code{path\_cost(path)} $=$ \code{len(path)} $-1$, and 0 for an empty path (\dec{04}) \\
\bottomrule
\end{tabularx}
\end{center}

\textbf{CSP (scheduling).} Variables are the task labels D1\dots Dn (all dirty cells, including
tasks with no constraint). The domain of each variable is the slots $\{1,\dots,T\}$, where $T$ comes
from line 1 of \code{Schedule.txt}; slots are integers and are shown as T1\dots T$T$ (\dec{06}).
There are three constraint types:
\begin{itemize}
  \item \emph{Precedence} $(D_i, D_j)$: $\mathrm{slot}(D_i) < \mathrm{slot}(D_j)$.
  \item \emph{Latest slot} $(D_i, T_k)$: $\mathrm{slot}(D_i) \le k$.
  \item \emph{Slot exclusivity} (implicit): no two tasks share a slot.
\end{itemize}

\textbf{Input formats.} \code{Environment.txt} has six lines: \code{rows,cols}; number of obstacles;
obstacle tuples; number of dirty cells; dirty-cell tuples; charger tuple. \code{Schedule.txt} has:
number of slots; number of precedence pairs; one \code{(Di,Dj)} per line; number of latest-slot
constraints; one \code{(Di,Tk)} per line. Both files are read at run time; nothing is hard-coded.
The parsers raise a single \code{InputError} naming the file, line and problem (\dec{07}). The rules
for whitespace, empty lists and value checks (duplicates, overlaps, out-of-bounds cells, unknown
labels) are \dec{08}--\dec{13}.

{\emergencystretch=3em Our own environment (\code{Environment.txt}) is a $10\times10$ room with 27 obstacles forming two
walls and a pocket, 5 dirty cells, and the charger at $(0,0)$. \code{Schedule.txt} has 6 slots, the
precedence chain D1$\to$D2$\to$D3$\to$D4, and the limits D4 $\le$ T4 and D2 $\le$ T3. D5 is
unconstrained, so it tests the ``every task gets a slot'' rule.\par}
